% ── Preámbulo del informe · LuaLaTeX + memoir ─────────────────────────────────────────────────────────────────────
\usepackage{iftex}
\ifluatex\else\errmessage{Este informe se compila con LuaLaTeX: latexmk -lualatex informe}\fi

% Idioma
\usepackage[english,spanish,es-noshorthands,es-tabla,es-nolayout]{babel}
\usepackage[autostyle=true]{csquotes}

% Tipografía
\usepackage{fontspec}
\usepackage{unicode-math}
\setmainfont{Libertinus Serif}[Numbers=OldStyle]
\setsansfont{Libertinus Sans}
\setmonofont{Source Code Pro}[Scale=0.82]
\setmathfont{Libertinus Math}
\newfontfamily\cifrasfont{Libertinus Serif}[Numbers={Lining,Monospaced}]
\newfontfamily\titularfont{Libertinus Sans}[Numbers=Lining]
\usepackage{microtype}

% Página
\setlrmarginsandblock{30mm}{28mm}{*}
\setulmarginsandblock{28mm}{30mm}{*}
\setheadfoot{14pt}{24pt}
\setheaderspaces{*}{10mm}{*}
\checkandfixthelayout
\setlength{\parindent}{1.2em}
\setlength{\parskip}{0pt}
\setlength{\emergencystretch}{2.5em}  % evita líneas que se salen por intervalos que no se cortan
\linespread{1.06}

% Colores: paleta categórica validada (dataviz, modo claro) y tintas de texto
% Sin la opción table: el colortbl instalado (2026) usa \insert@pcolumn, que el array de MiKTeX (2023) no define,
% y rompe toda columna p{} o X. El informe no usa filas coloreadas.
\usepackage{xcolor}
\definecolor{tinta}{HTML}{1B1B1A}
\definecolor{tinta2}{HTML}{52514E}
\definecolor{tinta3}{HTML}{8A8983}
\definecolor{reja}{HTML}{E4E3DE}
\definecolor{fondo}{HTML}{F6F5F1}
\definecolor{serieA}{HTML}{2A78D6}  % adaptado
\definecolor{serieB}{HTML}{EB6834}  % monitor de Tierra
\definecolor{serieC}{HTML}{1BAF7A}  % Tierra con ancho calibrado / solo días propios
\definecolor{azulnoche}{HTML}{0D366B}
\definecolor{azulmedio}{HTML}{1C5CAB}
\definecolor{azulclaro}{HTML}{CDE2FB}
\definecolor{alto}{HTML}{C0392B}
\definecolor{medio}{HTML}{B9770E}
\definecolor{bajo}{HTML}{5D6D7E}
\color{tinta}

% Matemática, unidades y tablas
\usepackage{mathtools}
% siunitx no se usa: la versión instalada no calza con el núcleo de LaTeX de MiKTeX. Las cifras llegan ya
% formateadas (coma decimal, espacio fino de miles) desde herramientas/construir_cifras.py.
\newcommand*{\un}[1]{\,\textup{#1}}
\usepackage{booktabs}
\usepackage{tabularx}
\usepackage{multirow}
\usepackage{longtable}
\setlength{\LTcapwidth}{\textwidth}  % mismo ancho de título que las demás tablas
\usepackage{threeparttable}
\usepackage{makecell}
\renewcommand{\arraystretch}{1.18}

% Listas
\usepackage{enumitem}
\setlist{itemsep=2pt,topsep=4pt,parsep=0pt}
\setlist[description]{font=\normalfont\sffamily\bfseries\color{azulnoche},leftmargin=1.4em,labelindent=0pt}

% Gráficos
\usepackage{graphicx}
\usepackage[section]{placeins}  % las figuras no se escapan de su sección
\setcounter{topnumber}{3}\setcounter{totalnumber}{4}
\renewcommand{\topfraction}{0.9}\renewcommand{\textfraction}{0.08}\renewcommand{\floatpagefraction}{0.75}
\usepackage{tikz}
\usetikzlibrary{babel,positioning,arrows.meta,calc,fit,backgrounds,shapes.geometric,decorations.pathreplacing,shadows}
\usepackage{pgfplots}
\pgfplotsset{compat=1.18}
\usepgfplotslibrary{groupplots,fillbetween}
\pgfplotsset{
  informe/.style={
    width=\linewidth, height=6cm,
    axis line style={tinta3}, tick style={tinta3},
    every tick label/.append style={font=\footnotesize\sffamily\color{tinta2}},
    label style={font=\footnotesize\sffamily\color{tinta2}},
    title style={font=\small\sffamily\bfseries\color{tinta}},
    legend style={font=\footnotesize\sffamily, draw=none, fill=none,
      /tikz/every even column/.append style={column sep=1.2em}},
    legend cell align=left,
    grid style={reja, line width=0.4pt},
    /pgf/number format/.cd, use comma, 1000 sep={\,},
  },
}

% Cajas
\usepackage[most]{tcolorbox}
\tcbset{
  base/.style={enhanced, breakable, boxrule=0pt, arc=1.5mm, left=4mm, right=4mm, top=2.5mm, bottom=2.5mm,
    fonttitle=\sffamily\bfseries\small, colback=fondo, before skip=8pt, after skip=10pt},
}
\newtcolorbox{clave}[1][]{base, unbreakable, colframe=serieA, borderline west={2.5pt}{0pt}{serieA},
  title={\color{serieA}Lo esencial}, coltitle=serieA, colbacktitle=fondo, attach title to upper={\quad}, #1}
\newtcolorbox{cuidado}[1][]{base, unbreakable, borderline west={2.5pt}{0pt}{serieB},
  title={\color{serieB}Cuidado}, colbacktitle=fondo, attach title to upper={\quad}, #1}
\newtcolorbox{decision}[1][]{base, unbreakable, borderline west={2.5pt}{0pt}{serieC},
  title={\color{serieC!80!black}Decisión}, colbacktitle=fondo, attach title to upper={\quad}, #1}
% hallazgo de la auditoría: {código}{gravedad}{color}{título}
\newtcolorbox{hallazgo}[4]{base, colback=white, boxrule=0.5pt, colframe=reja, borderline west={3pt}{0pt}{#3},
  title={\textcolor{#3}{#1\enspace·\enspace#2}\hfill\textcolor{tinta}{#4}}, coltitle=tinta, colbacktitle=white,
  toptitle=1mm, bottomtitle=0mm, titlerule=0pt}

% Código
\usepackage{listings}
\lstset{basicstyle=\ttfamily\small, columns=fullflexible, keepspaces=true, breaklines=true,
  backgroundcolor=\color{fondo}, frame=none, xleftmargin=3mm, xrightmargin=3mm, aboveskip=6pt, belowskip=6pt,
  literate={á}{{\'a}}1 {é}{{\'e}}1 {í}{{\'i}}1 {ó}{{\'o}}1 {ú}{{\'u}}1 {ñ}{{\~n}}1 {→}{{$\to$}}1}

% Bibliografía
\usepackage[backend=biber, style=authoryear-comp, maxcitenames=2, mincitenames=1, maxbibnames=8, minbibnames=6,
  uniquelist=false, uniquename=false, giveninits=true, dashed=false, backref=true, doi=true, url=false, isbn=false,
  eprint=true, sorting=nyt]{biblatex}
\addbibresource{generado/referencias.bib}
% La nota «Metadatos de …» que deja resolver_bib2.py dice de dónde salió la entrada; no va en la lista impresa
\DeclareSourcemap{\maps[datatype=bibtex]{\map{\step[fieldsource=note, match=\regexp{^Metadatos\s+de}, final]
  \step[fieldset=note, null]}}}
\DefineBibliographyStrings{spanish}{backrefpage={citado en la p\'agina}, backrefpages={citado en las p\'aginas}}
\renewcommand*{\bibfont}{\small}
\setlength{\bibitemsep}{3pt}

% Enlaces y referencias cruzadas
\usepackage[hidelinks=false, colorlinks=true, linkcolor=azulmedio, citecolor=azulmedio, urlcolor=azulmedio,
  pdfauthor={Benjamín Campos}, pdftitle={Transfer learning Tierra → espacio en NOVA},
  pdfsubject={Informe técnico}, pdfkeywords={transfer learning, ajuste fino, PatchTST, reposo en cama, NOVA},
  bookmarksnumbered=true]{hyperref}
\usepackage[spanish,noabbrev]{cleveref}
% Los nombres van en \extrasspanish: cleveref los repone cada vez que se vuelve al español (tras el abstract)
\newcommand{\nombresreferencias}{%
  \crefname{table}{tabla}{tablas}\Crefname{table}{Tabla}{Tablas}%
  \crefname{figure}{figura}{figuras}\Crefname{figure}{Figura}{Figuras}%
  \crefname{chapter}{capítulo}{capítulos}\Crefname{chapter}{Capítulo}{Capítulos}%
  \crefname{section}{sección}{secciones}\Crefname{section}{Sección}{Secciones}%
  \crefname{subsection}{sección}{secciones}\Crefname{subsection}{Sección}{Secciones}%
  \crefname{appendix}{apéndice}{apéndices}\Crefname{appendix}{Apéndice}{Apéndices}%
  \crefname{equation}{ecuación}{ecuaciones}\Crefname{equation}{Ecuación}{Ecuaciones}%
  \crefformat{equation}{ecuación~(##2##1##3)}\Crefformat{equation}{Ecuación~(##2##1##3)}%
}
\addto\extrasspanish{\nombresreferencias}
\AtBeginDocument{\nombresreferencias}

% Títulos de capítulo y sección
\makechapterstyle{nova}{%
  \renewcommand*{\chapterheadstart}{\vspace*{18mm}}
  \renewcommand*{\printchaptername}{}
  \renewcommand*{\chapternamenum}{}
  \renewcommand*{\printchapternum}{%
    {\titularfont\fontsize{48}{48}\selectfont\color{azulclaro}\thechapter}\par\vspace{-6mm}}
  \renewcommand*{\afterchapternum}{}
  \renewcommand*{\printchapternonum}{}
  \renewcommand*{\chaptitlefont}{\titularfont\bfseries\Huge\color{azulnoche}\raggedright}
  \renewcommand*{\printchaptertitle}[1]{\chaptitlefont ##1\par\vspace{3mm}{\color{serieA}\rule{24mm}{1.2pt}}}
  \renewcommand*{\afterchaptertitle}{\par\vspace{12mm}}
}
\chapterstyle{nova}
\setsecheadstyle{\titularfont\bfseries\Large\color{azulnoche}\raggedright}
\setsubsecheadstyle{\titularfont\bfseries\large\color{tinta}\raggedright}
\setsubsubsecheadstyle{\titularfont\bfseries\normalsize\color{tinta2}\raggedright}
\setbeforesecskip{-3.2ex plus -1ex minus -.2ex}
\setaftersecskip{1.6ex plus .2ex}
\setsecnumdepth{subsection}
\maxtocdepth{subsection}

% Cabeceras
\makepagestyle{nova}
\makeevenhead{nova}{\small\sffamily\color{tinta3}\leftmark}{}{\small\sffamily\color{tinta3}\thepage}
\makeoddhead{nova}{\small\sffamily\color{tinta3}\rightmark}{}{\small\sffamily\color{tinta3}\thepage}
\makeheadrule{nova}{\textwidth}{0pt}
\makepsmarks{nova}{%
  \createmark{chapter}{both}{shownumber}{}{.\ }%
  \createmark{section}{right}{shownumber}{}{.\ }}
\makepagestyle{plainnova}
\makeevenfoot{plainnova}{}{\small\sffamily\color{tinta3}\thepage}{}
\makeoddfoot{plainnova}{}{\small\sffamily\color{tinta3}\thepage}{}
\aliaspagestyle{chapter}{plainnova}
\pagestyle{nova}

% Leyendas
\captionnamefont{\small\sffamily\bfseries\color{azulnoche}}
\captiontitlefont{\small\color{tinta2}}
\captiondelim{.\enspace}
\setlength{\belowcaptionskip}{2pt}
\setfloatadjustment{figure}{\centering}
\setfloatadjustment{table}{\centering}

% Índices
\renewcommand*{\cftchapterfont}{\titularfont\bfseries}
\renewcommand*{\cftchapterpagefont}{\titularfont\bfseries}
\setlength{\cftbeforechapterskip}{6pt}

% Cifras generadas desde la salida del kernel (generado/cifras.tex). Una cifra que falta queda marcada y avisada.
\newcommand*{\cifra}[1]{%
  \ifcsname cifra@#1\endcsname\csname cifra@#1\endcsname
  \else\textbf{\textcolor{red}{??#1??}}\PackageWarning{informe}{Cifra indefinida: #1}\fi}
\newcommand*{\defcifra}[2]{\expandafter\def\csname cifra@#1\endcsname{#2}}

% Atajos
\newcommand*{\var}[1]{\texttt{#1}}
\newcommand*{\modelo}[1]{\textsf{#1}}
\newcommand*{\IC}[2]{\,[#1;\,#2]}
\newcommand*{\pct}{\,\%}
